\section{Statistical Model of CPC Decision-Making}\label{sec_methodology}

\subsection{Model}

To investigate the partial effects of case characteristics on case outcome, we developed and estimated a statistical model of CPC decision-making. We model each agenda item at a CPC hearing as having three possible outcomes:
\begin{enumerate}[start=0]
\item The project is denied or the decision is postponed;
\item The project is approved with conditions;
\item The project is approved without further conditions.
\end{enumerate}
Each project proposal $i$ is assumed to have a latent quality variable $y_i^\ast$ which determines the likelihood of the three outcomes. $y_i^\ast$ is modeled as a linear function of observed project characteristics and bureaucratic factors, $\mathbf{X}_i$, plus an error term $\epsilon_i$:
\begin{align}
y_i^\ast = \mathbf{X}_i \beta + \epsilon_i
\end{align}
The latent quality of the project proposal determines its outcome at the CPC hearing. Let $y_{i} \in \{0, 1, 2\}$ denote project $i$'s outcome. The relationship between $y_i^\ast$ and $y_i$ is as follows:
\begin{align}
y_i = \begin{cases}
0 \text{ if } y_i^\ast < \mu_0 \\
1 \text{ if } \mu_0 \leq y_i^\ast < \mu_1 \\
2 \text{ if } \mu_1 \leq y_i^\ast
\end{cases}
\end{align}
with $\mu_0 < \mu_1$. $\epsilon_{i}$ is assumed the follow a standard logistic distribution, and therefore the model is an ordered logit model, and the outcomes are monotonic in $\mathbf{X}_i \beta$. The parameters $\beta$, $\mu_0$, and $\mu_1$ are estimated by maximum likelihood.

\subsection{Estimation Results}

To operationalize the model, we include in $\mathbf{X}_i$ the following variables: indicators for whether the project is a residential, mixed-use, or non-residential development, the natural log of the square footage of the project, the height of the project in feet, the log-base-2 of the number of support letters, the log-base-2 of the number of opposition letters, the order in which the item appears on the agenda, the number of agenda items for that date's hearing, whether the item is on the consent calendar, and the case atypicality.\footnote{We chose not to include other physical characteristics such as number of dwelling units, number of affordable units, and parking spaces because they are highly correlated with the other physical dimensions. We also include dummy variables for whether the case is missing height or square footage information.} In addition, we include an increasingly comprehensive set of dummy variables, including dummies for the case suffix categories, the Council District to which the project belongs, the year in which the case was heard, and the text embeddings cluster the case belongs to.

Table \ref{tab_ologit_results} reports the estimated model coefficients for four different specifications. In column 1, the model includes only dummies for case suffix category, but not for Council District, year, or embedding cluster. In column 2, the model adds Council District dummies. In column 3, the model further adds year dummies. In column 4, the model adds embedding cluster dummies. The purpose of the successive inclusion of dummies is to show robustness of the coefficients of interest to the inclusion of further controls. Column 4 is the preferred specification.

Table \ref{tab_ologit_results} shows that although physical dimensions were associated with outcomes in the univariate analysis, the statistical association disappears when other factors are included in the regression. The coefficients on height and square footage are neither large nor statistically significant. Similarly, in the univariate analysis there was a statistically significant association between development type and outcomes. But in the multivariate regression analysis, the association between project type and outcome is no longer statistically significant.

The variables which are revealed to have a statistically significant association with outcomes are: the number of opposition letters, the order of the item in the agenda, whether the case is on the consent calendar, and the case atypicality. Using column 4 as the preferred specification, the results show that: i) a doubling in the number of opposition letters is associated with a \gn{NumOpposeCoef} reduction in log odds of full approval; ii) moving an agenda item one spot down in the agenda is associated with a \gn{AgendaOrderCoef} increase in log odds of full approval; iii) being on the consent calendar is associated with a \gn{ConCalCoef} increase in log odds of full approval; iv) and a one standard deviation increase in atypicality is associated with a \gn{AtypicalityCoef} decrease in log odds of full approval.

To translate these results into more interpretable numbers, we report the marginal effects of each variable in Table \ref{tab_ologit_marginal_effects}. The results indicate that for an ``average'' case: i) doubling the number of opposition letters is associated with a \gn{NumOpposeME2}\% reduction in the probability of full approval, but a \gn{NumOpposeME1}\% increase in the probability of approval with conditions, and a \gn{NumOpposeME0}\% increase in the probability of delay or denial; ii) moving an item down one spot in the agenda is associated with a \gn{AgendaOrderME2}\% increase in the probability of full approval, but a \gn{AgendaOrderME1}\% decrease in the probability of approval with conditions, and a \gn{AgendaOrderME0}\% decrease in the probability of delay or denial; iii) being on the consent calendar is associated with a \gn{ConCalME2}\% increase in the probability of full approval, and a \gn{ConCalME1}\% decrease in the probability of approval with conditions, and a \gn{ConCalME0}\% decrease in the probability of delay or denial; iv) a one standard deviation increase in atypicality is associated with a \gn{AtypicalityME2}\% decrease in the probability of full approval, but a \gn{AtypicalityME1}\% increase in the probability of being approved with conditions, and a \gn{AtypicalityME0}\% increase in the probability of being delayed or denied.

Overall, these results suggest that bureaucratic and political factors drive marginal decision-making at the CPC. If the CPC was primarily motivated by approving small residential developments and denying large non-residential ones, then we would expect to see statistically significant coefficients on project type and physical dimensions. Similarly, if the CPC was primarily motivated by strict adherence to the letter of the law in the zoning code, then we would not expect to find statistically significant influence of public sentiment and logistical issues, once planning characteristics are controlled for. Yet, that is what we find.

The signs of the coefficients are consistent with a theory of bureaucratic decision-making under time constraints, cognitive constraints, and oversight risk. Decisions on unfamiliar projects may be postponed because they require greater cognitive bandwidth, which is scarce under time constraints. Contested projects are also more likely to be postponed because they expose Commission members to greater political and reputational scrutiny. By contrast, routine and uncontested projects proceed with little resistance, as reflected by the consent calendar. 

\subsection{Bureaucratic capacity or latent quality?} \label{sec_latent_quality}

We have interpreted the effect of agenda order and atypicality as arising primarily from time constraints and cognitive constraints. That is, we have interpreted their effects as matters of bureaucratic capacity. A natural question that arises is whether their effects could instead be due to lower underlying quality.  If atypical proposals are of lower underlying quality, or if items earlier in the agenda are of lower underlying quality, then we can no longer interpret their association with delay as a matter of bureaucratic capacity, but rather as a matter of underyling quality issues with the proposal.

To test this, we ask whether delayed cases are more likely to be eventually denied or have conditions attached to them. If atypical and agenda order were merely proxies for latent quality, we would expect them to predict worse eventual outcomes regardless of whether the case was initially delayed; that is, they should be significantly more likely to be denied or approved with conditions in subsequent hearings. To test this hypothesis, we remove all delay decisions from the estimation sample and estimate a binary logit model in which the outcomes are 1 if the case is approved and 0 if the case is either approved with conditions or denied. The results for this regression are shown in Column 1 of Table \ref{tab_blogit_results}. In all four specifications, the coefficients on atypicality and agenda order becomes substantially attenuated. The coefficients on atypicality is no longer statistically significant in any of the four specifications, whereas the coefficients on agenda order are only statistically significant in the fourth regression, and only at the 10\% level. Thus, atypicality and agenda order have more statistical power in explaining delay than they do in explaining eventual outcome. This suggests that the association of agenda order and atypicality to delay cannot be attributed entirely to poor latent quality.

To further test our interpretation that the results are driven by bureaucratic capacity, we consider whether atypicality matters less after a case has already been postponed once and has come back to the CPC for another hearing. To do this, we limit the sample only to cases that had already been delayed once within our sample period.  We then estimate a binary logit model where the outcome is 1 if the case is approved or approved with conditions and 0 if the case is delayed again or denied. In the model, we interact atypicality with the number of times the case had been heard previously. Our hypothesis is that atypicality will matter less for cases already seen before, and that the effect of atypicality weakens the more times the case has been heard. The results are reported in columns 2-6 of Table \ref{tab_blogit_results}. Each column reports results with different combinations of controls.\footnote{We limit the number of controls in each specification due to the limited sample size of re-heard cases. Moreover, consent calendar is removed from the regressions because previously postponed cases are almost never placed on the consent calendar.} In each column, the coefficient on atypicality is statistically insignificant. The coefficient on the interaction term between atypicality and number of previous hearings is positive as hypothesized, but statistically insignificant in all specifications. If atypicality is merely a proxy for poor latent quality, we would expect atypicality to predict further delay and denial, even after the case has been heard multiple times. We therefore conclude that the effect of atypicality on delay is primarily driven by lack of familiarity, rather than latent proposal quality.

\subsection{Strategic agenda setting?}

Another issue we address is whether agenda item ordering or placement on the consent calendar may be politically motivated, or otherwise strategically determined. To explore this, we conduct two exercises. First, we estimate a binary logit model to uncover factors which predict placement on the consent calendar. Second, we estimate a OLS regression to find which factors are associated with the item's agenda order for non-consent-calendar items.\footnote{The consent calendar is always heard first before non-consent-calendar items.} 

Table \ref{tab_endogenous_agenda} shows the estimated coefficients for the above models. Column 1 shows the results for the consent calendar model. We find that public support and opposition are the most important determinants of placement on the consent calendar. Cases with a lot of support and opposition letters (more controversial cases) are less likely to be placed on the consent calendar. To test whether the CPC is responsive to political concerns beyond public opinion, we also include the number of support letters submitted by public officials. CPC members may be more responsive to elected officials if they have political ambitions themselves, or are otherwise concerned about their standing within political circles. However, we do not find evidence that letters by public officials are more influential than letters from the general public. We therefore find that public controversy is the main factor determining placement on the consent calendar. We do not find evidence for strategic placement beyond that intended purpose.

Column 2 of Table \ref{tab_endogenous_agenda} shows the results for the agenda order model. None of the factors appear to predict agenda order with statistical significance. We therefore conclude that there is no evidence suggesting strategic ordering of items in the agenda. Thus, the effect of agenda order on outcome that we estimated in Table \ref{tab_ologit_results} is more plausibly driven by bureaucratic or cognitive frictions, like decision fatigue or pressure to make faster decisions as time runs out.

\subsection{What happens to delayed cases?}

We now assess the extent to which delays at the CPC have a meaningful impact on economic outcomes. \citet{gabrielkung2025} showed that approval delays are an important driver of reduced housing production. They estimated that a 25\% reduction in approval times would have increased housing production in Los Angeles by 12.7\% between the years of 2010 and 2022. Thus, delays at the CPC could have a meaningful impact on real estate development if they add significant time to the overall approvals process. 

Table \ref{tab_delayed_outcomes} investigates what happens to the delayed cases in our dataset. There were \gn{NumDelayedCases} cases which were ever delayed in our dataset.\footnote{This number differs from the \gn{CasesContinued} instances of delay reported in Table \ref{tab_result_unanimity} because in Table \ref{tab_result_unanimity} counts multiple hearings of the same case as separate units.} Of these, \gn{NumResolved} eventually received a non-delay outcome within our dataset. The average time between first observed delay and final outcome was \gn{DaysToResolve} days. This is an underestimate because a) we don't observe the final outcome for some of the cases (right-censored resolution times) and b) we don't necessarily observe the first hearing (left-censored start times). But even this underestimate of \gn{DaysToResolve} days is economically significant: \citet{gabrielkung2025} report that the median approval time for multifamily housing development projects in Los Angeles from 2010 to 2022 was about 524 days, so \gn{DaysToResolve} days is a meaningful addition. Moreover, it is interesting to note that between the initial delay and final outcome, the majority of the delayed cases (\gn{ResolvedNoChangesPct}\%) did not change in terms of their requests, and some (\gn{ResolvedMinorChangesPct}\%) came back with only minor changes. Only \gn{ResolvedMajorChanges} of the delayed cases came back with any major changes to their requests; and even among these, only \gn{ResolvedMajorReduce} came back with reduced requests, while the other \gn{ResolvedMajorExpand} actually expanded their requests. The fact that a delay at the CPC does not usually cause the applicant to substantially modify their requests again suggests that the CPC's decision to delay is not necessarily correlated with the quality of the proposal, or else one would expect the developers to modify the proposal to increase the chance of approval.

\subsection{Policy implications}

We now discuss policy implications. Thus far, we have shown that public opposition and case atypicality predict delay at the CPC. Is there anything that can be done to reduce these frictions? 

To assess that question, we explored how different case characteristics moderated the effect of public support and opposition, agenda order, and atypicality. We found that the embedding cluster the case belongs to had the most significant moderating effects.\footnote{We also tried combinations of other variables, such as whether the case utilizes an incentive program or whether the case requires a conditional use permit. But the embedding clusters themselves had the most statistically significant effects. This suggests that the embedding clusters capture real information relevant to CPC decision-making, information that is not well captured by any single other case characteristic.} Recall that embedding cluster 0 captures large developments requiring tract map or subdivision approval, cluster 1 captures discretionary entitlements, conditional use permits, and code amendments, and cluster 2 consists primarily of residential projects utilizing incentives. Table \ref{tab_policy_imps} shows the results of ordered logit regressions in which public support and opposition, agenda order, consent calendar, and atypicality are interacted with embedding cluster dummies. Column 1 interacts embedding clusters with the number of support and opposition letters. Column 2 interacts the clusters with agenda order and consent calendar. Column 3 interacts the clusters with atypicality. Column 4 interacts the clusters with all of the above. 

We find evidence that opposition letters are most effective in embedding cluster 1, the group of cases requiring discretionary entitlements and conditional use permits. This makes sense because this is the group where the CPC members have most discretion. By comparison, opposition letters are less effective in cluster 2, which is mostly incentive program housing, because in order to qualify for incentive programs, these developments have already conformed to specific requirements in the planning code, thus removing discretion from CPC members. Indeed, many of the cases in cluster 2 are appeals of previous non-discretionary approvals, usually invoking CEQA (a state law) as the grounds for appeal. These appeals are mostly denied, and the project approved. These results suggest that introducing standardized pathways for approval, and thereby removing discretionary action from the approvals process, may be an effective policy for reducing the frictions caused by public opposition.

We do not find any evidence that the effects of agenda order or consent calendar are moderated by the embedding cluster. Indeed, there is little variation in whether an item is on the consent 



\subsection{Discussion}

%We have found evidence that bureaucratic and political factors are statistically significant drivers of decision-making at the CPC. We do not claim, however, that those factors matter more than the underlying quality of the project or its adherence to proper procedures and standards. It may simply be that by the time a project lands on the CPC's desk, it has already been pre-screened for quality, such that basic quality issues are no longer what determines approval at the margin. Unfortunately, we are not able to test this hypothesis with the data that we have, and we leave such work to future research. Our claim is much more modest: that decision-making at the CPC is slowed down by bureaucratic and political factors that may potentially be orthogonal to the overall planning objectives of the City. Thus, bureaucratic and political delay should register in the cost side of any cost-benefit analysis of new regulation. This is, of course, known quite generally, but we provide novel empirical evidence and some quantification of the effect.

Cognitive limits are inherently unobservable. They leave no entry in an agenda or a minute, and can only be inferred indirectly from patterns in observed decisions. Our estimates provide the basis for one such inference. If approval at the CPC hinged primarily on the fundamentals of a proposal, the fundamentals should carry more explanatory weight. Instead, project type and physical dimensions lose their association with outcomes once controls are included, while the variables describing the \emph{process}---public opposition, case atypicality, agenda order, and placement on the consent calendar---retain theirs.

The decision process described in Section~\ref{sec_background} offers an interpretation of this pattern, where the opposition penalty reflects oversight risk and the atypicality penalty comes from scarce attention. But an alternative reading is that these variables simply proxy for latent quality. Atypical cases may be atypical because something is wrong with them. Three features of the data resist this interpretation. First, once delayed cases are set aside, atypicality loses its power to predict final outcomes. Atypicality predicts whether a case is delayed, not how it is ultimately decided, supporting a processing friction rather than a quality defect. Second, the association between atypicality and delay is stronger when the case is re-heard by commissioners who had not seen it before. This is more consistent with lack of familiarity driving the result rather than a quality defect. Third, delayed cases overwhelmingly return unchanged: roughly three quarters come back with identical requests and are then resolved. If delay signaled a quality problem, applicants would presumably revise. Instead, for the most part, they simply wait.

Nonetheless, it is plausible that by the time a case reaches the CPC, it has already been screened for basic quality, so that the frictions we document operate only at the margin among pre-vetted proposals. We cannot test these possibilities with the data at hand, and we flag them as open questions rather than resolve them. Our claim is modest: decision-making at the CPC is slowed by bureaucratic and political factors that may be orthogonal to the overall planning objectives of the City. Bureaucratic and political delay should accordingly register on the cost side of any cost-benefit analysis of new regulation. 

